\documentclass[../../../include/open-logic-section]{subfiles}

\begin{document}
	
\olfileid{sth}{z}{nat}
\olsection{في اختيار الأعداد الطبيعية}

%لنسلّم بإمكان تبرير المبدأ غير الصوري~\stagesinf{}.
%لكن يجدر أيضًا التعليق على المسلّمة الخاصة التي استخرجناها من ذلك المبدأ غير الصوري.

حدّدنا في~\olref[infinity-again]{defnomega} مدلول «الأعداد الطبيعية»
تحديدًا صريحًا. وليس هذا الاصطلاح دعوى ميتافيزيقية، وإنما هو قرار بأن
\emph{نعامل} مجموعات مخصوصة معاملة الأعداد الطبيعية. وقد سبق ذكر بعض
دواعيه في
\olref[sfr][rel][ref]{sec}، و\olref[sfr][arith][ref]{sec}،
و\olref[sfr][infinite][dedekindsproof]{sec}؛ ويمكن الآن ضبطها أكثر.

مسلّمة اللانهاية التي أخذنا بها تتبع~\citet{VonNeumann1925}، وكان
يمكن أن نختار بدلها مسلّمة أخرى:

\begin{defish}
\emph{مسلّمة زرملو للانهاية في~\citeyear{Zermelo1908Untersuchungen}.}
توجد مجموعة~$A$ يكون فيها~$\emptyset \in A$ و
$(\forall x \in A)\{x\} \in A$.
\end{defish}

ولو أخذنا مسلّمة زرملو دون المسلّمة المستمدة من فون نويمان، لثبت كذلك
وجود مجموعة لا نهائية بمعنى ددكند، ثم جبر ددكندي. ويبقى العنصر المميز
في بناء زرملو~$\emptyset$، وهو نائبنا عن~$0$، لكن التطبيق الحقني يكون
$x \mapsto \{x\}$ بدل~$x \mapsto x \cup \{x\}$. وأظهر الفروق أن
زرملو يجعل~$2$ هو~$\{\{\emptyset\}\}$، ونجعله نحن، مع فون نويمان،
$\{\emptyset, \{\emptyset\}\}$.

فلم نرجح إحدى مسلّمتي اللانهاية؟ العلة العملية الكبرى أن نهج فون
نويمان يحسن «التوسع» إلى الأعداد المتناهية التجاوز، وسيبدأ بحث ذلك في
\olref[ordinals][]{chap}. أما مجرد \emph{إجراء الحساب} فيستوي فيه
النهجان. فإذا قيل إن الأعداد الطبيعية \emph{هي} مجموعات، وجب أن نسأل:
\emph{أي المجموعات؟}

وقد ذاع هذا السؤال الدقيق بفضل~\citet{Benacerraf1965}. وهو أشهر أفراد
ظاهرة مرت بنا مرارًا: فنظرية المجموعات تتيح \emph{محاكاة} أنواع من
الموجودات نفهمها حدسيًا، كالأعداد الحقيقية والنسبية والصحيحة والطبيعية،
وكالأزواج المرتبة والدوال والعلاقات؛ لكنها لا تعيّن قط محاكاة
\emph{وحيدة}. بل تبقى \emph{دائمًا} بدائل تؤدي الغرض نفسه.

\end{document}



